\documentclass[11pt,a4paper]{ctexart}
\usepackage[margin=18mm]{geometry}
\usepackage{xcolor}
\usepackage{graphicx}
\usepackage{fontspec}
\usepackage{amsmath}
\usepackage{unicode-math}
\usepackage[most]{tcolorbox}
\usepackage{enumitem}
\usepackage{titlesec}
\usepackage{needspace}
\setCJKmainfont{Songti SC}
\setCJKsansfont{Hiragino Sans GB}
\setmainfont{Avenir Next}
\setsansfont{Avenir Next}
\setmathfont{STIX Two Math}
\definecolor{ttblue}{HTML}{25508E}
\definecolor{ttmuted}{HTML}{5C6069}
\definecolor{ttlight}{HTML}{E8F0FC}
\definecolor{ttaccent}{HTML}{BE502D}
\definecolor{ttwarm}{HTML}{FFF6E8}
\titleformat{\section}{\Large\bfseries\sffamily\color{ttblue}}{}{0pt}{}
\titleformat{\subsection}{\large\bfseries\sffamily\color{ttblue}}{}{0pt}{}
\setlist[itemize]{leftmargin=1.5em,itemsep=0.22em,topsep=0.25em}
\XeTeXlinebreaklocale "zh"
\XeTeXlinebreakskip = 0pt plus 1pt
\emergencystretch=3em
\sloppy
\pagestyle{empty}
\newtcolorbox{chainbox}{colback=ttlight,colframe=ttblue!20,boxrule=0.6pt,arc=2mm,left=2mm,right=2mm,top=1mm,bottom=1mm}
\newtcolorbox{callout}[1]{colback=ttwarm,colframe=ttaccent!45,title={#1},fonttitle=\sffamily\bfseries\color{ttaccent},boxrule=0.7pt,arc=2mm,left=2.5mm,right=2.5mm,top=1.5mm,bottom=1.5mm}
\newcommand{\slideimage}[3]{%
\begin{center}
\includegraphics[page=#1,width=#2\linewidth]{#3}\\[-2pt]
{\footnotesize\color{ttmuted}原 PPT 第 #1 页}
\end{center}}
\begin{document}
{\sffamily\color{ttmuted}TTDS 10}\\[3pt]
{\Huge\sffamily\bfseries\color{ttblue}IR Evaluation 2: Test Collections、Pooling 与显著性检验}\\[6pt]
图文讲义版：用原 PPT 作为视觉锚点，按“概念故事线 + 直觉解释 + 考试速记”整理。\\[6pt]
\begin{chainbox}
\sffamily Cranfield paradigm 回顾 $\rightarrow$ test collection 的来源：web logs 或 evaluation campaigns $\rightarrow$ TREC tracks 与 task $\rightarrow$ topic 不等于 query：description 与 narrative 定义相关性 $\rightarrow$ relevance judgments 的成本问题 $\rightarrow$ pooling 用多个系统的 top results 构造可评估候选集 $\rightarrow$ Cohen's kappa 衡量标注者一致性 $\rightarrow$ web A/B testing 与 offline evaluation 互补 $\rightarrow$ significance testing 判断分数提升是否可靠
\end{chainbox}
\slideimage{1}{0.72}{/Users/huez/Documents/ttds/lec/ttds25_10evaluation2.pdf}
\begin{callout}{这节课一句话}
这节课一句话：TTDS 10 不是继续发明新指标，而是解释这些指标背后的“标准答案”从哪里来，以及平均分差异何时才算真实改进。
\end{callout}
\Needspace{0.38\textheight}
\section{1. 1. Test collection 从哪里来}
\slideimage{3}{0.62}{/Users/huez/Documents/ttds/lec/ttds25_10evaluation2.pdf}
\begin{callout}{人话直觉}
评价系统需要“考卷”和“答案”。Google 这类公司可以从真实用户行为里造题，学术社区则常靠大型评测活动一起造题。
\end{callout}
\noindent\textbf{精确定义 / 机制：} Web search companies 可以监控 traffic、click/session logs，并对 selected queries 做人工标注；非 web search 任务通常没有现成大规模日志，所以需要 evaluation campaigns 构建 collection、topics 和 judgments。第 3 页强调 test collection 很贵，campaigns 的价值就是把成本集中投入，形成可复用资源。

\noindent\textbf{考试问法：} 若问 why evaluation campaigns，答构建标准 test collections 成本高，需要社区共享任务、数据和 judgments 来比较系统。

\Needspace{0.38\textheight}
\section{2. 2. TREC 与 track：不是一个比赛，而是一套任务生态}
\slideimage{4}{0.62}{/Users/huez/Documents/ttds/lec/ttds25_10evaluation2.pdf}
\begin{callout}{人话直觉}
TREC 像 IR 的标准化考场；不同 track 是不同科目的试卷，比如 medical、legal、microblog、CLIR。
\end{callout}
\noindent\textbf{精确定义 / 机制：} TREC 是由 NIST 赞助的主要 IR evaluation campaign。一个 track 围绕特定文档类型、领域或任务组织搜索任务，提供 collection、topics、judgments 和评估协议。它的核心价值不是让某个系统赢一次，而是让不同方法在同一任务定义下可比较、可复现、可累积。

\noindent\textbf{考试问法：} 若问 TREC track 是什么，答特定 domain/genre/task 的评测任务集合，并给 Medical/Legal/Microblog/CLIR 例子。

\Needspace{0.38\textheight}
\section{3. 3. Topic vs Query：考试最容易混淆的一对}
\slideimage{5}{0.62}{/Users/huez/Documents/ttds/lec/ttds25_10evaluation2.pdf}
\begin{callout}{人话直觉}
query 是用户打在搜索框里的几个字；topic 是给 assessor 的完整案件说明书，告诉他什么算相关、什么不算。
\end{callout}
\noindent\textbf{精确定义 / 机制：} TREC topic 通常包含 num、title/query、description 和 narrative。title 是短 query；description 解释信息需求；narrative 明确 relevance criteria，减少 assessor 之间的理解差异。第 5 页的 Health and Computer Terminals 示例展示了同一短标题背后还需要细化的判定边界。

\noindent\textbf{考试问法：} 若问 topic 为什么比 query 更适合评测，答它降低歧义，让人工 judgments 更一致、更可复现。

\Needspace{0.38\textheight}
\section{4. 4. Relevance judgments 为什么不能 exhaustive}
\slideimage{6}{0.62}{/Users/huez/Documents/ttds/lec/ttds25_10evaluation2.pdf}
\begin{callout}{人话直觉}
如果 50 个 topics、百万文档，每个 topic-document pair 都人工判断，就像让老师批五千万份卷子，现实上不可行。
\end{callout}
\noindent\textbf{精确定义 / 机制：} Exhaustive assessment 的规模是 topics \(\times\) documents，成本巨大；random sampling 又因为 relevant documents 稀少，可能抽不到足够相关样本。因此需要一种聪明的候选收集方法：只判断最可能被系统返回、也最可能影响排名指标的文档。这正是 pooling 的动机。

\noindent\textbf{考试问法：} 若问为什么不用 random sampling，答 relevant docs sparse，随机抽样效率低且可能无法覆盖真正重要的相关文档。

\Needspace{0.38\textheight}
\section{5. 5. Pooling：用多个系统帮 assessor 找候选答案}
\slideimage{6}{0.62}{/Users/huez/Documents/ttds/lec/ttds25_10evaluation2.pdf}
\begin{callout}{人话直觉}
Pooling 像让很多学生先交出他们觉得可能正确的答案，老师只批这些高频候选；不是批完整宇宙，但覆盖了最可能影响比较的区域。
\end{callout}
\noindent\textbf{精确定义 / 机制：} 典型流程是：系统每个 topic 提交 top 1000，取每个系统 top 100 进入 pool，去重并随机排序，由 topic creator 判断；未评估文档通常视为 irrelevant，再计算 MAP 等指标到 rank 1000。Pooling 要有效，需要大量 reasonable systems，而且系统不能全做同一件事，否则 pool 会偏向单一方法。

\noindent\textbf{考试问法：} 若考 pooling，必须写步骤、未评估文档处理方式、为什么需要 many diverse systems，以及它不是 exhaustive 但在实践中可用。

\Needspace{0.38\textheight}
\section{6. 6. Cohen's kappa：相关性判断也要测可靠性}
\slideimage{7}{0.62}{/Users/huez/Documents/ttds/lec/ttds25_10evaluation2.pdf}
\begin{callout}{人话直觉}
两个 assessor 同意很多次不一定厉害；如果数据极不平衡，他们可能靠碰巧也常同意。Kappa 要扣掉 chance agreement。
\end{callout}
\noindent\textbf{精确定义 / 机制：} Cohen's kappa 衡量两个 judges 的 agreement 超过随机一致性的程度。P(A) 是 observed agreement，P(E) 是 by chance 的 expected agreement。\(\kappa\)=1 表示完全一致，\(\kappa\)=0 表示只达到随机水平，\(\kappa\)<0 比随机还差；第 8 页给出经验阈值：>0.8 good，0.67–0.8 fair，<0.67 评测基础可疑。

\[
\kappa = \frac{P(A)-P(E)}{1-P(E)}
\]
\noindent\textbf{考试问法：} 若问公式，解释每个项；若问 relevance judgments 的问题，答主观性 + multiple assessors + agreement measure。

\Needspace{0.38\textheight}
\section{7. 7. Web Search A/B Testing：真实用户流量里的实验}
\slideimage{9}{0.62}{/Users/huez/Documents/ttds/lec/ttds25_10evaluation2.pdf}
\begin{callout}{人话直觉}
Offline evaluation 像模拟考；A/B testing 像把新菜单给 1\% 真顾客试吃，直接看他们会不会更满意。
\end{callout}
\noindent\textbf{精确定义 / 机制：} A/B testing 的目的通常是测试单个 innovation。大多数用户继续用 old system，少量流量导向 new system，用 clickthrough on first result 等 automatic measures 观察用户满意度是否改善。它需要已有大型 search engine 和足够流量，因此是大公司最信任但不是所有任务都能使用的方法。

\noindent\textbf{考试问法：} 若问 web search 公司如何评估，答 traffic、logs、selected query labeling、A/B testing；若问为什么 recall 难测，说明 web 全体 relevant docs 不可知。

\Needspace{0.38\textheight}
\section{8. 8. Significance Testing：平均分高不是结案陈词}
\slideimage{12}{0.62}{/Users/huez/Documents/ttds/lec/ttds25_10evaluation2.pdf}
\begin{callout}{人话直觉}
系统 B 平均 MAP 高一点，可能是真强，也可能只是这组 queries 运气好。显著性检验是在问：这个差异像不像随机波动？
\end{callout}
\noindent\textbf{精确定义 / 机制：} Null hypothesis 通常是假设两个系统没有真实差异，observed differences 只是随机波动。Significance test 试图拒绝 null hypothesis；test power 是正确拒绝 null 的概率，增加 query 数量通常提高 power。IR 中常用 two-tailed t-test，\(\alpha\)=0.05；如果不能假设 per-query differences 近似正态，可以考虑 Wilcoxon。

\noindent\textbf{考试问法：} 若问为什么平均分更高还不够，答需要统计显著性；若问 t-test 假设，答 per-query effectiveness differences 来自 normal distribution 且均值差为 0 是 null。

\section{考试速记版}
\begin{itemize}
\item Test collection = documents + topics + relevance judgments + measure。
\item Topic = query/title + description + narrative；query 只是短文本。
\item Pooling = 多系统 top results 合并去重后人工判断。
\item Pooling 需要 many reasonable and diverse systems。
\item Unevaluated docs 通常当作 irrelevant。
\item Cohen's kappa 扣除 chance agreement 后衡量 assessor agreement。
\item A/B testing 适合已有大规模真实流量的 web search。
\item Significance test 用来判断平均分提升是否可能只是随机波动。
\end{itemize}
\begin{callout}{一段稳妥考场回答}
A reusable IR test collection is normally built from a document collection, a set of topics, relevance judgments and an evaluation measure. In TREC-style evaluation, a topic is richer than a query because it includes a short title/query, a description of the information need, and a narrative specifying relevance criteria. Since exhaustive judging is too expensive and random sampling rarely finds enough relevant documents, pooling is used: multiple diverse systems submit top-ranked documents, the top results are merged and de-duplicated, assessors judge the pool, and unjudged documents are usually treated as non-relevant. Finally, a higher average score is not sufficient evidence that one system is better; a significance test, often a two-tailed t-test at \(\alpha\)=0.05 or a non-parametric alternative such as Wilcoxon, is needed to assess whether the improvement is likely to be real.
\end{callout}
\end{document}
